\documentclass[12pt]{article}
\usepackage[utf8]{inputenc}
\usepackage[T1]{fontenc}
\usepackage{hyperref}
\title{Uma' Lasan Language Resources for K-12 Curriculum Development}
\author{Compiled by Claude Sonnet 5.5 (claude-sonnet-5-5) for Andr\'as Kornai}
\date{}
\begin{document}
\maketitle

\section{Introduction}
Uma' Lasan (ISO 639-3 xky) is an Austronesian language of the Kenyah group, spoken in Sarawak (Malaysia) and in Indonesian Borneo. Related names include Western Kenyah and Uma Baka. It is distinct from Uma' Lung (ulu), which has its own code. One thesis places a community at Long Pari village, Bulungan regency. I found no published dictionary, grammar or school materials in my searches.

\section{Primary Educational Resources}
None found. I found no primer or textbook, and no school programme.

\section{Linguistic and Reference Materials}
\begin{itemize}
\item Wikipedia, ``Uma' Lasan language'': an overview. \url{https://en.wikipedia.org/wiki/Uma%CA%BC_Lasan_language}
\item Wikipedia, ``Kenyah languages'': context on the wider group. \url{https://en.wikipedia.org/wiki/Kenyah_languages}
\item A thesis on affixes in the Uma' Lasan language, based on speech recorded at Long Pari. It is circulated on Scribd. I could not verify its publication status. \url{https://www.scribd.com/document/285477182/ABSTRACT-docx}
\item Glottolog reference record for Soriente 1997. I did not verify its relevance to Uma' Lasan. \url{https://glottolog.org/resource/reference/id/95925}
\end{itemize}

\section{Religious and Community Resources}
\begin{itemize}
\item Joshua Project profile, ``Kenyah, Western in Malaysia'': it lists Uma' Lasan as the primary language and reports that Scripture translation has started. \url{https://joshuaproject.net/people_groups/print/15372/MY}
\end{itemize}

\section{Digital and Audio Resources}
\begin{itemize}
\item Global Recordings Network, Uma' Lasan: audio and audio-visual Bible lessons, with downloadable MP3 and MP4 packages. \url{https://globalrecordings.net/en/language/xky}
\item Global Recordings Network, Uma' Lasan (Uma Baka variety). \url{https://glorecnet-website.k.globalrecordings.net/en/language/11863}
\end{itemize}

\section{Curriculum Development Context}
With about 2,800 speakers, and with not all young people speaking it, the language has a real chance of being used in education. The documented base is thin: there are recordings and one affix study, and no dictionary or grammar. A first step would be a lexicon and orthography built with the community, probably drawing on the Kenyah literature, in Malaysian and Indonesian settings.
\end{document}